\subsection{Envelope criteria and extension across the degeneracy set}
\label{subsec:semipositive-envelope-extension}

The capacity criterion above asks for a genuine smaller-background obstacle envelope. Two sufficient conditions come from different sources: invariant bounded mass controls the contact measures of stretched truncations, while a sublinear singularity bound supplies explicit competitors. Both conditions give zero Lelong numbers. Under dual attainment, a globally smooth root of rank at least two then permits smooth extension across the analytic degeneracy set.

We retain the connected compact manifold $X$ of dimension $n\geq2$, the Hermitian background $\omega$ and the smooth positive probability volume $\mathrm{d}V$. All obstacle inequalities concern canonical potentials. We will use one consequence of \cref{lem:gauduchon-uniform-lelong}: if, for every $0<\delta<1$, a genuine $v_\delta\in\PSH(X,\delta\omega)$ satisfies $v_\delta\leq u$, then
\[
 0\leq\nu(u,x)\leq\nu(v_\delta,x)
       =\delta\nu(v_\delta/\delta,x)\leq\delta C_\omega.
\]
Thus $u$ has zero Lelong numbers. The bound is uniform in $\delta$ because $v_\delta/\delta$ has the fixed background $\omega$.

\begin{theorem}\label{thm:invariant-mass-threshold-envelopes}\index{Envelope!invariant mass criterion}
Suppose $\ddc(\omega^j)=0$ for $1\leq j<n$. Put $V=\int_X\omega^n$. For every nonzero $0\leq f\in L^1(\mathrm{d}V)$, with $M=\int_Xf\,\mathrm{d}V$, one has
\[
 b_\omega(f)=\frac VM.
\]
Every threshold subsolution satisfies $\mathrm{MA}^{\mathrm{ext}}_\omega(u)=(V/M)f\,\mathrm{d}V$ and has mass $V$. Its obstacle envelope $P_{\delta\omega}(u)$ is genuine for every $0<\delta<1$, and $u$ has zero Lelong numbers. More generally, these last two conclusions hold for every $\omega$-plurisubharmonic potential of extended mass $V$.
\end{theorem}
\begin{proof}
By \cref{thm:mass-invariance-characterization}, every bounded $\omega$-plurisubharmonic potential has mass $V$. Truncation therefore bounds every extended mass above by $V$. For $f_k=\min(f,k)$, integration of the bounded equation gives $c_\omega(f_k)=V/\int_Xf_k\,\mathrm{d}V$. The definition \eqref{eq:individual-truncated-normalizer} gives $b_\omega(f)=V/M$. A threshold subsolution has mass at least $b_\omega(f)M=V$. Its mass is consequently $V$, and the nonnegative difference from $b_\omega(f)f\,\mathrm{d}V$ has zero mass.

Let now $u$ have extended mass $V$. Normalize $u\leq0$, put $\mu=\mathrm{MA}^{\mathrm{ext}}_\omega(u)$, and fix $a>1$. For $u_j=\max(u,-j)$, let $w_j=P_\omega(au_j)$. Constants give $-aj\leq w_j\leq0$. The bounded contact theorem \cref{thm:bounded-contact-support} concentrates its measure on $C_j=\{w_j=au_j\}$. Canonical almost-everywhere order gives $w_j\leq au_j$ everywhere. Both potentials are $a\omega$-plurisubharmonic, so positivity and ordered contact give
\[
 \mathrm{MA}_\omega(w_j)
 \leq\mathbf1_{C_j}\mathrm{MA}_{a\omega}(w_j)
 \leq a^n\mathbf1_{C_j}\mathrm{MA}_\omega(u_j).
\]
For integers $j\geq N$, locality and the bounded mass identity yield
\[
 \int_{\{u_j\leq-N\}}\mathrm{MA}_\omega(u_j)
       =V-\mu(\{u>-N\})=\mu(\{u\leq-N\}).
\]
Choose $N$ with $a^n\mu(\{u\leq-N\})<V$. This is possible because $\mu$ is finite and gives no mass to the pole set. If $\sup_Xw_j<-aN$, its contact set lies in $\{u_j<-N\}$, contradicting the preceding inequalities and $\int_X\mathrm{MA}_\omega(w_j)=V$. Hence $\sup_Xw_j\geq-aN$ for $j\geq N$; for smaller $j$ the bound follows from $w_j\geq-aj$.

The envelopes decrease. Their uniform lower supremum prevents collapse to $-\infty$, by normalized quasi-plurisubharmonic compactness. Their genuine limit $w$ satisfies $w\leq au$, including at poles. With $\delta=1/a$, the function $\delta w$ is a genuine $\delta\omega$-plurisubharmonic competitor below $u$. Thus $P_{\delta\omega}(u)$ is genuine at every scale. The Lelong estimate preceding the theorem completes the proof. An additive normalization does not affect any conclusion.
\end{proof}

\begin{proposition}\label{prop:sublinear-threshold-envelope}\index{Envelope!sublinear singularity criterion}
Suppose a genuine quasi-plurisubharmonic $\psi\leq0$ satisfies $\ddc\psi\geq-C\omega$ for some $C>0$, and a genuine $\omega$-plurisubharmonic $u$ satisfies the canonical inequality
\[
 u\geq-A(1-\psi)^\alpha-B,
 \qquad A>0,\quad B\in\mathbb R,\quad0<\alpha<1.
\]
Then every $P_{\delta\omega}(u)$ is genuine, for $0<\delta<1$, and $u$ has zero Lelong numbers.
\end{proposition}
\begin{proof}
Choose $\alpha<\beta<1$ and $\varepsilon>0$ with $\varepsilon\beta C\leq\delta$. The number
\[
 K=B+\sup_{t\geq1}(At^\alpha-\varepsilon t^\beta)
\]
is finite, and $v_\delta=-\varepsilon(1-\psi)^\beta-K\leq u$. This composition is increasing and convex, with first derivative at most $\varepsilon\beta$ and nonnegative second derivative. Apply the quasi-plurisubharmonic chain inequality to $\max(\psi,-j)$ and pass to the decreasing limit. It gives $\ddc v_\delta\geq-\varepsilon\beta C\omega\geq-\delta\omega$. The limit is genuine because $\psi$ is finite almost everywhere. Its canonical pole values are $-\infty$, where the comparison still holds. Thus $v_\delta$ is an actual competitor at each scale. The preceding uniform Lelong estimate proves the last assertion.
\end{proof}

\subsubsection{A globally smooth root of rank at least two}

\begin{theorem}\label{thm:rank-two-dual-smooth-extension}\index{Gauduchon duality!smooth extension across degeneration}
Assume the semipositive attained setting of \cref{thm:semipositive-dual-exact-equation}, with greatest normalized threshold solution $\Phi_f$ and analytic degeneracy set $Z$. Suppose the maximizing form has the representation $T=\gamma^{n-1}$ for a globally smooth semipositive $(1,1)$-form $\gamma$ whose complex rank is at least two everywhere. If $\Phi_f$ has zero Lelong numbers, then it is smooth on $X$. This applies under either sufficient condition above. Every threshold subsolution is then smooth, since these subsolutions differ by constants.
\end{theorem}
\begin{proof}
Put $U=X\setminus Z$ and $A=\omega+\ddc\Phi_f$. The existing attained theorem makes $A$ smooth on $U$ and gives its exact global equation. The root $\gamma$ is positive on $U$. Stokes gives $\int_XA\wedge T=1$, while attainment gives
\[
 \int_U(A^n)^{1/n}(\gamma^n)^{(n-1)/n}
       =b_\omega(f)^{1/n}D_{f^{1/n}}(T)=1.
\]
The arithmetic--geometric mean gap is nonnegative and continuous on $U$, and its integral must vanish. Equality yields $A=r\gamma$ there for smooth $r\geq0$. Where $A^n=0$, equality forces $A=0$, so the same formula holds with $r=\operatorname{tr}_\gamma A/n=0$.

Since $\ddc\Phi_f$ is closed,
\[
 \mathrm dr\wedge\gamma+r\,\mathrm d\gamma=\mathrm d\omega
       \quad\hbox{on }U.
\]
For a semipositive form of rank at least two, the real-linear map $W_\gamma:\xi\mapsto\xi\wedge\gamma$ from one-forms to three-forms is injective. In a diagonal real unitary frame, for each coordinate of $\xi$ choose a positive coordinate two-plane not containing it. Its triple coefficient is that coordinate times the positive eigenvalue; no other disjoint diagonal pair contributes. Vanishing of the wedge therefore forces every coordinate to vanish.

Choose an auxiliary metric. The left inverse $L_\gamma=(W_\gamma^*W_\gamma)^{-1}W_\gamma^*$ is smooth and uniformly bounded on compact $X$. Thus
\begin{equation}\label{eq:rank-two-extension-first-order}
 \mathrm dr=\alpha-r\beta,\qquad
 \alpha=L_\gamma(\mathrm d\omega),\quad
 \beta=L_\gamma(\mathrm d\gamma),
\end{equation}
with smooth one-forms on $X$. In particular $|\mathrm d\log(1+r)|\leq C$ on $U$.

The analytic set permits short avoiding paths. In a coordinate ball, stratify it into countably many smooth strata of real dimension at most $2n-2$. For a fixed outside point, endpoints of segments meeting one stratum lie in a cone of dimension at most $2n-1$, a Lebesgue-null set. For two nearby outside points choose a generic intermediate point in a ball of radius comparable to their distance. Both connecting segments avoid $Z$ and their combined length is bounded by a fixed multiple of that distance. A finite connected coordinate cover, with outside points chosen in its overlaps, also gives a global uniform bound for connecting lengths in $U$.

The logarithmic gradient bound now bounds $r$ on $U$, and \eqref{eq:rank-two-extension-first-order} bounds its derivative. The local short paths give a unique locally Lipschitz extension across $Z$. Since $Z$ is Lebesgue-null, its weak derivative satisfies the same equation on $X$. The right side is continuous, so $r$ is $C^1$; iteration makes it smooth.

Set $S=r\gamma$. It is smooth semipositive and $\mathrm dS=\mathrm d\omega$ everywhere by density. The current $R=A-S$ is closed and supported on $Z$. It is positive: it equals the positive coefficient-measure part of $A$ on $Z$, since $A=S$ off $Z$ and $S$ has no mass there. The analytic support theorem recalled in \cref{sec:psh-divisorial-estimates} expresses $R$ as a nonnegative sum of divisorial integration currents; components of codimension at least two carry no such current. Locally $S-\omega$ is smooth and closed and has a smooth $\ddc$-potential. The divisorial coefficients are therefore the Lelong numbers of $\Phi_f$, and all vanish. Consequently $\ddc\Phi_f=S-\omega$ is smooth. Taking a Euclidean trace and applying elliptic regularity gives a smooth potential; its canonical representative agrees with it everywhere.
\end{proof}

The smooth root and its rank bound are explicit hypotheses. A smooth semipositive maximizing $(n-1,n-1)$-form need not have such a root across its degeneracy set. On a surface the rank bound already forces $\gamma>0$ everywhere.

\subsubsection{Two attained instances}

\begin{example}\label{ex:hopf-surface-invariant-attainment}\index{Gauduchon duality!Hopf surface attainment}
On the Hopf surface $X=(\mathbb C^2\setminus\{0\})/\langle z\mapsto2z\rangle$, the background $\omega=|z|^{-2}\ddc|z|^2$ has invariant bounded mass. It admits nonconstant smooth threshold solutions with actual smooth positive maximizing tests.
\end{example}
\begin{proof}
The form descends under $z\mapsto2z$ and is positive. For $s=|z|^2$ and $F(s)=s^{-1}$, the identity $sF''+2F'=0$ gives $\ddc\omega=\ddc(s^{-1})\wedge\ddc s=0$. It is nonclosed. Choose a nonconstant smooth real $q$ and small nonzero $\varepsilon$ such that $A=\omega+\ddc(\varepsilon q)>0$. Put $V=\int_X\omega^2$, $\mathrm{d}V=\omega^2/V$, $f=A^2/\mathrm{d}V$ and $u=\varepsilon q-\sup_X\varepsilon q$. Invariant mass gives $\int_Xf\,\mathrm{d}V=V$, hence $b_\omega(f)=1$ by \cref{thm:invariant-mass-threshold-envelopes}.

The test $T=A/V$ satisfies $\ddc T=0$, $\int_X\omega\wedge T=1$ and $D_{\sqrt f}(T)=1$. For every normalized smooth positive Gauduchon test $Q$, arithmetic--geometric mean and Stokes give $D_{\sqrt f}(Q)\leq\int_XA\wedge Q=1$. Thus $T$ attains the supremum. The attained uniqueness theorem identifies $\Phi_f=u$. This is a bounded nonconstant instance.
\end{proof}

\begin{example}\label{ex:quadric-rank-two-dual-attainment}\index{Gauduchon duality!quadric ramification}
For $n\geq3$, a smooth quadric double cover of $\mathbb P^n$ has an actual semipositive maximizing test whose smooth root has rank at least two everywhere and degenerates along the ramification divisor.
\end{example}
\begin{proof}
Let
\[
 X=\{[w:z_0:\cdots:z_n]:w^2=\textstyle\sum_{i=0}^nz_i^2\}
       \subset\mathbb P^{n+1},\qquad F([w:z])=[z].
\]
The quadratic is nonsingular and the map is finite of degree two. Its smooth ramification divisor is $Z=\{w=0\}$. Local branch coordinates show that $\mathrm dF$ has rank $n-1$ on $Z$ and $n$ elsewhere. Let $\omega$ be the restricted normalized Fubini--Study form and $A=F^*\omega_{\mathrm{FS}}$. Then
\[
 u=\log\frac{\sum_i|z_i|^2}{|w|^2+\sum_i|z_i|^2}
\]
is global and smooth, $-\log2\leq u\leq0$, $\sup_Xu=0$, and $A=\omega+\ddc u$. Branched change of variables gives $\int_XA^n=2$, and Stokes gives $V=\int_X\omega^n=2$.

Set $\mathrm{d}V=\omega^n/V$, $f=A^n/\mathrm{d}V$, $T=A^{n-1}/V$ and $\gamma=A/V^{1/(n-1)}$. The closed background gives $b_\omega(f)=1$. Stokes gives $\ddc T=0$ and $\int_X\omega\wedge T=1$; the determinant weights give $D_{f^{1/n}}(T)=1$. Arithmetic--geometric mean, valid also for semipositive $A$, bounds every competing normalized positive test by $\int_XA\wedge Q=1$. Thus $T$ is an actual maximizer and $\Phi_f=u$ by attained uniqueness. Its root has rank $n-1\geq2$ on $Z$ and is positive elsewhere, as required.
\end{proof}

The invariant-mass and sublinear conditions answer restricted instances of the envelope question. General moment or entropy bounds have not been shown to supply either condition or the smooth-root rank hypothesis.
